\documentclass{ximera}

\begin{document}
	\author{Stitz-Zeager}
	\xmtitle{Answers for App Derivatives}{}

\mfpicnumber{1} \opengraphsfile{ExercisesforAppDerivatives} % mfpic settings added 

\begin{enumerate}

\item increasing:  $(-\infty, -1)$, $(2, \infty)$;  decreasing:  $(-1,2)$;  local max:  $(-1,8)$;  local min:  $(2, -19)$.

\smallskip

\item increasing:  $(-1,1)$;  decreasing:  $(-\infty, -1)$, $(1, \infty)$;  local max:  $(1,5)$;  local min:  $(-1, -5)$.

\smallskip

\item increasing:  $\left( \frac{3}{2}, \infty\right)$; decreasing:  $\left( -\infty, \frac{3}{2}\right)$;  local (absolute) min:  $\left(\frac{3}{2}, -\frac{3}{2 \sqrt[3]{2}}\right)$

\smallskip

\setcounter{HW}{\value{enumi}}
\end{enumerate}

\begin{enumerate}
\setcounter{enumi}{\value{HW}}


\item  concave up: $\left(\frac{1}{2}, \infty\right)$;  concave down: $\left( - \infty, \frac{1}{2} \right)$; inflection point:  $\left(\frac{1}{2}, \frac{11}{2}\right)$

\smallskip

\item concave up:  $\left( -\sqrt{3}, 0 \right)$, $\left( \sqrt{3}, \infty \right)$;  concave down:    $\left(- \infty,  -\sqrt{3} \right)$, $\left(0,  \sqrt{3} \right)$;  inflection points:  $\left( -\sqrt{3}, -\frac{5 \sqrt{3}}{2} \right)$, $(0,0)$, $\left( \sqrt{3}, \frac{5 \sqrt{3}}{2} \right)$
\smallskip

\item  concave up:  $(-\infty, 2)$, $(3, \infty)$;  concave down:  $(2,3)$;  inflection points: $(2,0)$, $(3, 3)$. 

\smallskip

\setcounter{HW}{\value{enumi}}
\end{enumerate}

\begin{enumerate}
\setcounter{enumi}{\value{HW}}

\item  \begin{enumerate}  \item To find the $x$-intercept, we set $ax+b = 0$  and get $x = -\frac{b}{a}$ provided $a \neq 0$.

\item  Solving  $f''(x) = 6ax + 2b = 0$, we get $x = -\frac{2b}{6a} = - \frac{b}{3a}$, provided $a \neq 0$.  

\smallskip

The graph of $f''(x) =  6ax + 2b$ is a line so we know on one side of $x= - \frac{b}{3a}$, $f''(x) > 0$ and on the other side, $f''(x) < 0$.  

\smallskip

Hence, the graph of the original function $y = f(x)$ changes concavity at $x = -\frac{b}{3a}$.

\end{enumerate}

\item\begin{enumerate}  \item  To solve $\overline{C}'(x) = \frac{x \, C'(x) - C(x)}{x^2} = 0$, we set the numerator,  $x \, C'(x) - C(x) = 0$. We get  $x \, C'(x)  = C(x)$ so $C'(x) = \frac{C(x)}{x} = \overline{C}(x)$.

\smallskip

\item  Divide both numerator and denominator of  $\overline{C}''(x) = \frac{x^2 \, C''(x) - 2x\, C'(x) + 2C(x)}{x^3}$ by $x$: 

\smallskip

$\overline{C}''(x) = \frac{x \, C''(x) - 2 \, C'(x) + 2\frac{C(x)}{x}}{x^2}$ and substitute  $\frac{C(x)}{x} = \overline{C}(x)$.

\smallskip

\item  If  $C'(x) = \overline{C}(x)$, then:

\smallskip

 $\overline{C}''(x)  = \frac{x\, C''(x) - 2\, C'(x) + 2\overline{C}(x)}{x^2} =  \frac{x\, C''(x) - 2\overline{C} (x) + 2\overline{C}(x)}{x^2}= \frac{x\, C''(x)}{x^2} =  \frac{C''(x)}{x}$.

\smallskip

\item  When $C'(x) = \overline{C}(x)$, we have that $\overline{C}'(x) = 0$ and $\overline{C}''(x) > 0$.  By the Second Derivative Test for Local Extrema,  Theorem \ref{secondderivatvetest}, the average cost $\overline{C}(x)$ has a minimum when $C'(x) = \overline{C}(x)$ 

\smallskip



\end{enumerate}

\newpage

\item \begin{enumerate}  \item

\begin{multicols}{2}

 $f(x) = 0$:  $x = 0, \pi$
 
  $f'(x) = 0$: $x = 2.029$ ($f_{+}'(0) = 0$.,  too.)

\end{multicols}

\smallskip

\item  \begin{multicols}{2}

 $f(x) > 0$: $(0, \pi)$
 
  $f(x) < 0$: none

\end{multicols}

\smallskip

\begin{multicols}{2}

 $f'(x) > 0$: $(0, 2.029)$
 
  $f'(x) < 0$: $(2.029, \pi)$

\end{multicols}

\smallskip

\begin{multicols}{2}

 $f''(x) > 0$: $(0, 1.077)$
 
  $f''(x) < 0$: $(1.077, \pi)$.

\end{multicols}

\end{enumerate}

\item Answers vary.  below is a sketch of the key features that should be included:


\begin{center}

\centerline{\includegraphics[width = 5in]{./AppDerivativesGraphics/graphfromderivativeexerciseanswer.png}}
\end{center}


\setcounter{HW}{\value{enumi}}
\end{enumerate}

\end{document}
